% The kappa-interval question: measured obstruction and architecture.
% Input at the end of the margins section.

\subsection{From points to intervals}
\label{sec:interval}

The natural strengthening of Theorem~\ref{thm:main} replaces each
margin by a $\kappa$-interval, so that $\alpha_\star(\kappa)$ becomes
an unconditional bound on a continuum.  For
Conditions~\ref{cond:fp}--\ref{cond:lc} this is what the strip
certificates already provide: those inequalities hold with margins of
order $10^{-1}$, and enclosing every input over a $\kappa$-slab costs
far less than that.  Condition~\ref{cond:2var} is different in kind,
and it is worth recording precisely why, because the obstruction is
quantitative and, we believe, intrinsic.

Certifying $S_\star\le 0$ over a slab means every certificate must
hold with the located parameters $(\alpha, q_0, \psi_0)$ replaced by
their hulls across the slab.  Three measurements fix the scales.
First, the hull widths cannot be made small: the located enclosures
carry irreducible floors (the inward fixed-point test resolves no
finer than the quadrature fuzz, about $3.5\cdot10^{-4}$ in $q_0$, and
$\alpha_\star$ itself drifts at $|d\alpha_\star/d\kappa|\approx
1$--$3.4$ across any slab), which translate into a moment-plane
origin ball of radius at least a few times $10^{-3}$ however narrow
the slab.  Second, the evaluation fuzz of every certificate grows
linearly with that ball: rerunning the $\kappa=-0.65$ sweep with
slab-hulled parameters (ball $2.3\cdot10^{-2}$) inflates cell
enclosures to $\pm 1.3$, and the scaling back to the floor ball
leaves $\pm 0.2$.  Third, the certificates that matter are thin:
the sweep's weakest certified cells at point margins sit at
$-5\cdot10^{-4}$, and Region~I's thinnest bands at $-4.7\cdot10^{-5}$,
because the two-variable function is genuinely flat at its degenerate
maximizer.  A $10^{-4}$-margin certificate cannot survive
$10^{-1}$-scale fuzz, at any slab width, with this evaluator
architecture.

The same comparison explains why the strip conditions posed no such
problem and points at the correct construction.  Near the maximizer
the certified Hessian gives quadratic control whose constants vary
slowly in $\kappa$ (the fixed-tilt determinant moves by about
$2\cdot10^{-4}$ per unit margin), so a $\kappa$-uniform neighborhood
estimate is available from point certificates plus a crude
$\kappa$-derivative bound --- crude is enough, because a Lipschitz
constant needs no thin margin.  Far from the maximizer the
certificates are fat and absorb hull fuzz directly.  What remains is
the intermediate annulus, where quadratic control must be extended
beyond the star radius with explicit remainder bounds before the two
regimes meet.  We record this as the open engineering of the margin
program: the obstruction is measured, the two easy regimes are in
hand, and the annulus estimate is the one genuinely new ingredient an
interval theorem requires.

The certified artifacts also price that ingredient.  Across the three
positive margins the Region~I band populations have nearly identical
margin profiles: at most one or two bands per margin certify below
$3\cdot10^{-4}$, at most seven below $3\cdot10^{-3}$, and more than
$95$ percent above $3\cdot10^{-2}$.  If each band carried a certified
bound $L_i$ on the $\kappa$-derivative of its certificate value, a
band would hold over a $\kappa$-radius $m_i/L_i$ of its margin
$m_i$, and re-certifying each band at its own pitch would cover
$10^{-2}$ of margin axis for $1.1$ times the cost of one point run
at unit derivative bound, $1.3$ times at a conservative bound of
three --- against roughly $3\cdot10^{5}$ band evaluations for
uniform slabs at the thinnest band's pitch.  The
compute is therefore not the obstruction; what is missing is the
derivative enclosure itself, one certified $\kappa$-derivative per
certificate family, with the located branch differentiated
implicitly.  The constant itself is measurable from the existing
artifacts: the band lattice is shared across margins up to the wedge
rotation, and $1292$ rectangles coincide exactly between the
$\kappa=0.05$ and $\kappa=0.0995$ manifests, giving certified
finite differences of each certificate value along the located
branch.  The measured sensitivities are strongly in the plan's
favor: the median band moves at $|\Delta S/\Delta\kappa|\approx 6$,
but the thin-trough bands move at $0.1$--$0.3$, so the thinnest band
(margin $5.6\cdot10^{-4}$) already holds over a $\kappa$-radius
near $2\cdot10^{-3}$, and with per-band measured pitches the count
is $1.04$ runs per $10^{-2}$ of margin axis --- about eight
Region~I run-equivalents for all of $[0.05,0.13]$.  A difference
over a $0.05$ gap bounds the mean drift rather than the local
derivative, so these numbers size the project rather than prove it;
the machinery an interval theorem still requires is the certified
per-certificate $\kappa$-derivative, with the same measurement then
repeated for the sweep, union, and star-interior families.
